% Source: https://artofproblemsolving.com/community/c3h1926988_selfsimilarity
% Stable identity: aops-topic-1926988-q01
For $c\ge0$ define $a_1=\sqrt c$ and $a_{n+1}=\sqrt{c+a_n}$ for $n\ge1$, and let $L(c)=\displaystyle\lim_{n\to\infty}a_n$ be the value of the infinite expression $\sqrt{c+\sqrt{c+\sqrt{c+\cdots}}}$.
\begin{enumerate}
\item[(a)] Find $L(0)$. The ``self-similarity'' method sets $x=\sqrt x$ and gets $x=0$ or $x=1$. Which is the value, and why is the other one extraneous?
\item[(b)] For $c>0$, prove that $L(c)=\dfrac{1+\sqrt{1+4c}}{2}$.
\item[(c)] Show that $L(c)\to1$ as $c\to0^+$, although $L(0)=0$.
\end{enumerate}
